To test the evaluation protocol beyond recurrent clustering and its gate, we
apply a fixed PELT plus segmented VAR-Lasso adapter to Study C's 30 paired
repeats at both durations and the original $300/150$ grid. PELT uses the same
raw window features, a penalised unknown segment count, and no K-means or
gate. Each detected contiguous interval is independently refit on its
nonoverlapping time range. This is a different segmentation/refit structure,
while retaining the local linear score input; it is our graph adapter to the
external change-point algorithm, not an off-the-shelf causal method.

Table~\ref{tab:external-framework} shows why a pooled improvement alone is
insufficient here too. PELT plus VAR raises pooled AUROC relative to raw by
$+0.013$ $[0.006,0.020]$ at $D=300$ and $+0.008$ $[0.001,0.014]$ at
$D=600$. State-AUROC differences are $+0.005$ $[-0.020,0.027]$ and
$-0.006$ $[-0.020,0.008]$: neither establishes a state-ranking gain or
equivalence. At $D=600$, its positive-change AP contrast is $-0.125$
$[-0.205,-0.044]$, whereas standard AP gives $-0.011$ $[-0.067,0.048]$.
The metric and zero-group distinctions therefore remain consequential
without the prototype gate.

All 60 adapter cases converge and retain 644 segment refits. PELT chooses
9--11 segments on 23 windows, often near the minimum-length limit; the
configuration is not retuned after observing this behaviour. The comparison
demonstrates the protocol on one additional structure, not broad algorithmic
generality. Exact settings, segment counts, retained failures policy and
paired pointwise intervals are in Supplement Section~S20.
